\documentclass[11pt,a4paper]{article}
% Common preamble for the hint sheets (\input from each file)
\usepackage[utf8]{inputenc}
\usepackage[T1]{fontenc}
\usepackage{lmodern}
\usepackage{amsmath,amssymb,amsthm,mathtools}
\usepackage[a4paper,margin=2.3cm]{geometry}
\usepackage{enumitem}
\usepackage{xcolor}
\usepackage[hidelinks]{hyperref}
\renewcommand{\leq}{\leqslant}
\renewcommand{\le}{\leqslant}
\renewcommand{\geq}{\geqslant}
\renewcommand{\ge}{\geqslant}
\setlength{\parindent}{0pt}
\setlength{\parskip}{0.35em}
% \hintstitle{Sheet title}{Class work / Homework}
\newcommand{\hintstitle}[2]{\begin{center}{\LARGE\bfseries Hints: #1\par}\vspace{0.5em}{\large #2\par}\end{center}
\noindent\textit{The hints to each problem are cumulative. Read one, try again, and open the next only if you are still stuck.}\par\medskip\hrule\medskip}
\newcommand{\problem}[1]{\par\bigskip{\large\bfseries Problem #1}\par\nopagebreak\smallskip}
\newcommand{\subpart}[1]{\par\smallskip\textbf{(#1)}\par\nopagebreak}
\newcommand{\hint}[1]{\par\smallskip\noindent\textcolor{black!70}{\textbf{Hint #1.}}\enspace}
\newcommand{\note}[1]{\par\smallskip\noindent\textit{Note.} #1\par}
\providecommand{\lt}{<}
\providecommand{\gt}{>}
\providecommand{\tightlist}{\setlength{\itemsep}{0pt}\setlength{\parskip}{0pt}}
\emergencystretch=3em

\begin{document}
\hintstitle{Integrals -- 2. Symmetry, the king's rule and periodicity}{Homework}

\problem{1}
\hint{1}
Substitute \(x\mapsto -x\): the integral becomes \(\int_{-a}^a\frac{f(x)}{1+e^{-x}}dx=\int_{-a}^a\frac{e^xf(x)}{1+e^x}dx\). Add it to the original: you get \(\int_{-a}^af=2\int_0^af\).
\hint{2}
\(\int_{-1}^1\frac{x^2}{1+e^x}dx=\int_0^1x^2dx=\frac13\) and \(\int_{-\pi/2}^{\pi/2}\frac{\cos x}{1+e^x}dx=\int_0^{\pi/2}\cos x\,dx=1\).

\problem{2}
\hint{1}
Pair the terms: \(\frac1{1+\tan x}+\frac1{1+\cot x}\), \(\frac1{1+\cos x}+\frac1{1+\sec x}\), \(\frac1{1+\sin x}+\frac1{1+\csc x}\). Each pair has the form \(\frac1{1+u}+\frac1{1+1/u}\).
\hint{2}
\(\frac1{1+u}+\frac{u}{u+1}=1\), so the integrand equals \(3\) wherever it is defined, and the integral is \(3x+C\).

\problem{3}
\hint{1}
The shift \(x\mapsto x+\pi\) sends \(\sin x\mapsto-\sin x\). So \(\int_0^{2\pi}\lfloor 2023\sin x\rfloor dx=\int_0^\pi\left(\lfloor y\rfloor+\lfloor -y\rfloor\right)dx\) with \(y=2023\sin x\).
\hint{2}
\(\lfloor y\rfloor+\lfloor -y\rfloor=-1\) unless \(y\) is an integer, which happens only at finitely many \(x\). So the answer is \(-\pi\).

\problem{4}
\hint{1}
Compute the first few nested expressions separately for \(x\ge0\) and \(x\lt0\). For the first one, \(g_1=\lvert x\rvert\), \(g_{k+1}=\lvert x+g_k\rvert\): for \(x\ge0\), \(g_k=kx\); for \(x\lt0\) the values alternate \(-x,0,-x,0,\dots\).
\hint{2}
With \(2026\) (even) occurrences, \(g_{2026}=0\) for \(x\lt0\) and \(2026x\) for \(x\ge0\), so the first integral is \(\int_0^12026x\,dx=1013\).
\hint{3}
For the second, the function is even; for \(x\ge0\) let \(h_0=x\), \(h_k=\lvert h_{k-1}-1\rvert\). By induction, on \([0,n]\) the function \(h_n\) is a zigzag between \(0\) and \(1\) with slopes \(\pm1\) (a triangle wave), so its integral over \([0,n]\) is \(\frac n2\). Hence the answer is \(2\cdot\frac{2023}2=2023\).

\problem{5}
\hint{1}
In the first integral \(x^{2024}dx\) is, up to the factor \(\frac1{2025}\), exactly \(d(x^{2025})\); so \(u=x^{2025}\) (or \(u=1-x^{2025}\), which maps \([0,1]\) to itself reversed) reduces it to a single power.
\hint{2}
First: \(\frac1{2025}\int_0^1u^{2025}du=\frac1{2025\cdot2026}\). Second: \(u=1-x\) gives \(\int(u^{2021}-u^{2020})\,du\) up to sign, so the answer is \(\frac{(1-x)^{2022}}{2022}-\frac{(1-x)^{2021}}{2021}+C\).

\problem{6}
\hint{1}
\(\cos x+\cos(x+\frac{2\pi}3)+\cos(x-\frac{2\pi}3)=\operatorname{Re}\left(e^{ix}(1+\omega+\omega^2)\right)\) with \(\omega=e^{2\pi i/3}\).
\hint{2}
\(1+\omega+\omega^2=0\), so the integrand is identically \(0\) and the integral is a constant \(C\).

\problem{7}
\hint{1}
Apply the king's rule to \(I=\int_a^b\frac{f}{f+g}\): since \(f(a+b-x)=g(x)\) and hence \(g(a+b-x)=f(x)\), it becomes \(\int_a^b\frac{g}{g+f}\). Add the two expressions.
\hint{2}
\(2I=b-a\). A classical example (Putnam 1987): \(\int_2^4\frac{\sqrt{\ln(9-x)}}{\sqrt{\ln(9-x)}+\sqrt{\ln(x+3)}}dx=1\), here \(a+b=6\) and \(f(x)=\sqrt{\ln(9-x)}\), \(g(x)=\sqrt{\ln(x+3)}\); no elementary antiderivative is known (and none is expected).

\problem{8}
\hint{1}
Simplify each summand for non-integer \(x\) (integers form a null set): \(\lfloor\lceil x\rceil\rfloor=\lceil x\rceil\), \(\lceil\lfloor x\rfloor\rceil=\lfloor x\rfloor\), \(\lfloor\{x\}\rfloor=0\), \(\{\lfloor x\rfloor\}=0\), \(\lceil\{x\}\rceil=1\), \(\{\lceil x\rceil\}=0\).
\hint{2}
The integrand is \(\lfloor x\rfloor+\lceil x\rceil+1=2\lfloor x\rfloor+2\) almost everywhere, so the integral is \(2\sum_{n=0}^{999}n+2000=1001000\).

\problem{9}
\hint{1}
\(\int_0^N\binom{\lfloor x\rfloor}{k}dx=\sum_{n=0}^{N-1}\binom nk\), and the hockey-stick identity gives \(\sum_{n=0}^{N-1}\binom nk=\binom N{k+1}\).
\hint{2}
So the answers are \(\binom{20}3=1140\) and \(\binom{2026}4=699938073450\).

\problem{10}
\hint{1}
Either note that \(F(\sin x)\), with \(F'=f\), is an antiderivative and is \(2\pi\)-periodic, or use \(x\mapsto\pi-x\) (together with periodicity): \(\sin\) is unchanged and \(\cos\) changes sign.
\hint{2}
General criterion: if a substitution \(x\mapsto c-x\) maps the interval onto itself, or \(x\mapsto x+T\) maps it onto itself modulo a period, and turns the integrand \(h\) into \(-h\), then \(I=-I\), so \(I=0\). Good test cases from the class sheet: 2(b) with \(x\mapsto\pi-x\), 2(c) with \(x\mapsto2\pi-x\), 2(d) with \(x\mapsto\pi-x\), and the odd integrands in Problem 4.
\end{document}
